% === CH04-THE-LAUNCHPAD ===
% Storymakers Method Report — Chapter 4
% "The 20% That Determines the Other 80%"

\section{The 20\% That Determines the Other 80\%}

\pullquote{The most important work in any presentation happens before a single slide is opened. It is invisible, unglamorous, and non-negotiable.}{Storymakers Principle}

\sourcefig{infographics/ch04-launchpad.jpg}{The Launchpad: SCQA, Big Idea, Three Pillars, and Audience Profile — the strategic foundation built before any slide is created.}

\subsection{Before the Slides: The Launchpad}

There is a moment in every presentation project where the outcome is effectively decided. It is not the moment the designer chooses a color palette. It is not the moment the data analyst selects the chart type. It is not the moment the speaker rehearses delivery. It is earlier --- much earlier --- in the quiet, uncomfortable phase where the presenter must answer a set of questions that most people would rather avoid.

We call this phase the Launchpad.\endnote{The Launchpad concept synthesizes frameworks from Barbara Minto's Pyramid Principle (1987), Nancy Duarte's audience-centric approach in ``Resonate'' (2010), and McKinsey's standard engagement start-up process. The term and specific integration are original to the Storymakers Method.}

The Launchpad is not a template. It is not a checklist to rush through. It is the strategic foundation upon which every subsequent decision rests --- the argument structure, the narrative arc, the visual hierarchy, the slide count, the level of detail. Get the Launchpad right, and the presentation almost builds itself. Get it wrong, and no amount of design polish or rehearsal will save it.

\begin{thesisbox}
\textbf{Core Thesis.} The first 20\% of effort invested in a presentation --- the strategic thinking before any slides exist --- determines the quality of the remaining 80\%. Most presenters invert this ratio, spending 80\% of their time on slides and 20\% (or less) on strategic foundations. This inversion is the single most reliable predictor of presentation failure.
\end{thesisbox}

The Launchpad has four components, each building on the previous one: (1) the SCQA Framework, defining the argument; (2) the Big Idea, crystallizing the core message; (3) the Three Pillars, structuring the evidence; and (4) the Audience Profile and Temperature Scale, calibrating the delivery.

\subsection{SCQA: The Argument Engine}

The SCQA framework originates with Barbara Minto, who developed it at McKinsey \& Company in the late 1960s as part of the Pyramid Principle.\endnote{Barbara Minto, ``The Pyramid Principle: Logic in Writing and Thinking,'' Financial Times / Prentice Hall, 3rd edition, 2009. Minto developed the original framework while working as McKinsey's first female MBA hire. The SCQA structure has since become the de facto standard for structured business argumentation across the consulting industry.} It is deceptively simple: four letters, four questions, and a rigorous logical chain that forces the presenter to confront the purpose of the communication before touching a single slide.

\textbf{S --- Situation.} What does the audience already accept as true? This is the common ground, the shared reality, the uncontroversial starting point. If the audience disagrees with your Situation, the rest of the framework collapses.

\textbf{C --- Complication.} What has changed? What tension, threat, or opportunity disrupts the Situation? The Complication is the engine of the argument --- it creates urgency. Without a Complication, there is no reason to present at all.

\textbf{Q --- Question.} What does the audience now need to resolve? The Question emerges naturally from the collision between Situation and Complication. It should be the question the audience is already forming in their minds.

\textbf{A --- Answer.} What is your response? This is the Big Idea --- the single, definitive answer to the Question. It is your thesis, your recommendation, your point of view.

\begin{center}
\begin{tabularx}{\textwidth}{>{\raggedright}p{2.5cm} >{\raggedright}p{4.5cm} >{\raggedright\arraybackslash}p{6.5cm}}
\toprule
\rowcolor{tableheader} \textcolor{white}{\textbf{Element}} & \textcolor{white}{\textbf{Function}} & \textcolor{white}{\textbf{Test Question}} \\
\midrule
Situation & Establish common ground & ``Would the audience nod along without objection?'' \\
\rowcolor{altrow}
Complication & Create urgency and tension & ``Does this make the status quo feel untenable?'' \\
Question & Channel the audience's thinking & ``Is this the question they are already forming?'' \\
\rowcolor{altrow}
Answer & Deliver the Big Idea & ``Can I say this in one sentence under 20 words?'' \\
\bottomrule
\end{tabularx}
\end{center}

\begin{insightbox}[SCQA IS NOT OPTIONAL]
SCQA is not an academic exercise. It is a diagnostic tool. If you cannot articulate each element clearly, your presentation has a structural defect that will manifest as audience confusion, disengagement, or --- worst of all --- the dreaded ``so what?'' at the end of your talk. Every failed presentation we have analyzed can be traced to a missing or malformed SCQA.
\end{insightbox}

\subsection{The Big Idea: One Sentence to Rule Them All}

The Answer in SCQA is not a paragraph. It is not a summary. It is one sentence --- what Nancy Duarte calls the ``Big Idea.''\endnote{Nancy Duarte, ``Resonate: Present Visual Stories that Transform Audiences,'' Wiley, 2010. Duarte defines the Big Idea as a single sentence that articulates ``your unique point of view'' and ``what's at stake'' for the audience. The constraint of a single sentence is deliberate: it forces radical clarity.}

The Big Idea has two mandatory components:

\begin{itemize}
  \item \textbf{A Point of View:} Not a neutral observation but a position. Not ``we should consider options'' but ``we must do X.''
  \item \textbf{What's at Stake:} The consequence of action or inaction. Not ``this is important'' but ``if we fail to act, we lose \$50M in market share within 18 months.''
\end{itemize}

The constraint is severe and intentional: the Big Idea must be expressible in fewer than 20 words. If it takes more than 20 words, it is either two ideas (split it), too vague (sharpen it), or not yet understood (think harder).

\begin{alertbox}[COMMON FAILURE]
\textbf{The ``Topic as Big Idea'' Trap.} The most common mistake is confusing a topic with an idea. ``Digital transformation'' is a topic. ``Our current systems will cost us 30\% of our customer base within two years unless we migrate by Q3'' is a Big Idea. Topics invite exploration. Big Ideas demand decision.
\end{alertbox}

\begin{center}
\begin{tabularx}{\textwidth}{>{\raggedright}p{5.5cm} >{\raggedright}p{5.5cm} >{\centering\arraybackslash}p{2.2cm}}
\toprule
\rowcolor{tableheader} \textcolor{white}{\textbf{Weak (Topic / Vague)}} & \textcolor{white}{\textbf{Strong (Big Idea)}} & \textcolor{white}{\textbf{Words}} \\
\midrule
``We need to talk about our Q3 results'' & ``Q3 revenue fell 18\%; pivoting to enterprise recovers margin by Q1'' & 13 \\
\rowcolor{altrow}
``An update on the product roadmap'' & ``Delaying v3 past June risks losing first-mover advantage in APAC'' & 12 \\
``Customer satisfaction trends'' & ``NPS dropped 22 points because onboarding takes 3x longer than competitors''' & 13 \\
\rowcolor{altrow}
``Thoughts on the reorg'' & ``Consolidating three BUs into one saves \EUR{12}M and cuts decision latency 60\%'' & 13 \\
``Some ideas about AI adoption'' & ``AI triage in support cuts resolution time 40\% and pays back in 5 months'' & 14 \\
\bottomrule
\end{tabularx}
\end{center}

Every strong Big Idea contains a specific claim, a measurable stake, and an implicit call to action. Every weak version is a topic label --- it tells the audience what the presentation is \emph{about} but not what it \emph{argues}.

\pullquote{Your Big Idea should make someone in the room slightly uncomfortable. If everyone nods in comfortable agreement, you have stated a truism, not an idea.}{Storymakers Principle}

\subsection{The Three Pillars: MECE Sub-Arguments}

A Big Idea without supporting evidence is an assertion. To become an argument, it needs structure. The Storymakers Method requires exactly three supporting pillars --- not two, not five, three --- each of which must satisfy the MECE principle: Mutually Exclusive, Collectively Exhaustive.\endnote{The three-pillar constraint draws on cognitive load research by George Miller (``The Magical Number Seven, Plus or Minus Two,'' 1956) and subsequent refinements by Nelson Cowan (``The Magical Number 4 in Short-Term Memory,'' 2001). Three pillars sit comfortably within working memory limits while providing sufficient structural depth.}

\textbf{Mutually Exclusive:} No pillar overlaps with another. If your three pillars are ``Revenue Growth,'' ``Sales Performance,'' and ``Commercial Results,'' you have one pillar wearing three hats.

\textbf{Collectively Exhaustive:} The three pillars together must fully support the Big Idea. If a skeptical audience member asks ``but what about X?'' and X is not covered by any pillar, the structure is incomplete.

Each pillar should:

\begin{itemize}
  \item \textbf{Have its own governing headline} --- a single sentence that states what this section proves.
  \item \textbf{Be supported by 2--4 pieces of evidence} --- data points, case studies, expert testimony.
  \item \textbf{Stand alone} --- if you presented only this pillar, it would still make a coherent sub-argument.
  \item \textbf{Connect back to the Big Idea} --- every pillar must visibly support the overarching thesis.
\end{itemize}

The result is Minto's Pyramid Principle in its purest application: the Big Idea at the apex, three pillars beneath it, and a layer of evidence beneath each pillar.\endnote{Gene Zelazny, ``Say It with Presentations,'' McGraw-Hill, 2006. Zelazny, McKinsey's longtime Director of Visual Communication, codified the pyramid-to-slides translation process that the Storymakers Method extends.}

\begin{diagnosisbox}
\textbf{Why Three?} Three is not arbitrary. It is the minimum number required for structural stability (two feels thin, four starts to blur) and the maximum that audiences reliably retain from a single sitting. McKinsey's internal presentation guidelines have required three key messages since the 1970s. Duarte's research found that TED talks with three supporting arguments scored significantly higher in audience retention than those with four or more.\endnote{Internal McKinsey communication guidelines, referenced in Ethan Rasiel, ``The McKinsey Way,'' McGraw-Hill, 1999. The cognitive basis for the rule of three is rooted in pattern recognition research; see also Maria Konnikova, ``Mastermind,'' Viking, 2013.}
\end{diagnosisbox}

\subsection{The Temperature Scale: Calibrating Intent}

Not every presentation has the same purpose. The Temperature Scale defines five levels that determine the rhetorical strategy, the evidence density, and the emotional register.\endnote{The Temperature Scale builds on Aristotle's three modes of persuasion (logos, ethos, pathos) and their relative weight depending on communication objective. The five-level formalization is original to the Storymakers Method.}

\begin{center}
\begin{tabularx}{\textwidth}{>{\centering}p{1cm} >{\raggedright}p{2.2cm} >{\raggedright}p{3.8cm} >{\raggedright}p{3cm} >{\raggedright\arraybackslash}p{3cm}}
\toprule
\rowcolor{tableheader} \textcolor{white}{\textbf{\#}} & \textcolor{white}{\textbf{Level}} & \textcolor{white}{\textbf{Objective}} & \textcolor{white}{\textbf{Dominant Mode}} & \textcolor{white}{\textbf{Example}} \\
\midrule
1 & Inform & Transfer knowledge, no action & Logos & Monthly dashboard review \\
\rowcolor{altrow}
2 & Align & Build shared understanding & Logos + Ethos & Strategy offsite kickoff \\
3 & Recommend & Propose a specific course of action & Logos + Ethos & Budget reallocation proposal \\
\rowcolor{altrow}
4 & Persuade & Overcome objections, change minds & All three modes & Board approval for acquisition \\
5 & Inspire & Ignite emotional commitment & Pathos + Ethos & Company transformation launch \\
\bottomrule
\end{tabularx}
\end{center}

\begin{keybox}[TEMPERATURE DETERMINES EVERYTHING]
The temperature level dictates slide design, narrative structure, evidence density, and delivery style. An ``Inform'' presentation that uses ``Inspire'' rhetoric feels manipulative. An ``Inspire'' presentation loaded with ``Inform'' data feels lifeless. Mismatched temperature is one of the most common --- and most invisible --- sources of audience disconnection. Most presenters default to Level 1 (Inform) regardless of the actual need. They dump data, walk through charts, and wonder why the audience did not act.\endnote{Nancy Duarte, ``How to Present to Senior Executives,'' Harvard Business Review, October 2012. Duarte notes that executives consistently complain about presentations that ``present data without a point of view.''}
\end{keybox}

\subsection{The Audience Profile: Who Are You Talking To?}

The final Launchpad component is the most neglected. Audience profiling is not demographics. It is a rigorous assessment of four dimensions:

\begin{enumerate}
  \item \textbf{What they already know.} Overestimate this, and you lose them to confusion. Underestimate it, and you insult their intelligence.
  \item \textbf{What they fear.} Every audience has anxieties. A presentation that ignores these fears feels tone-deaf. One that acknowledges them builds trust.
  \item \textbf{What they can decide.} Presenting a \$50M investment case to a team that controls a \$2M budget is a structural mismatch, not a persuasion failure.
  \item \textbf{What they value.} Speed versus thoroughness. Innovation versus stability. Revenue versus margin. These values determine which evidence lands.\endnote{Audience analysis frameworks draw on stakeholder mapping methodologies and persuasion psychology. See: Garr Reynolds, ``Presentation Zen,'' New Riders, 2012, Chapter 2, on audience-centric design.}
\end{enumerate}

\begin{alertbox}[THE AUDIENCE MISMATCH]
\textbf{Presenting to the wrong audience is worse than presenting badly.} A technically perfect presentation delivered to people who cannot act on it is organizational waste. Before building a single slide, confirm: Does this audience have the knowledge, the authority, and the motivation to make the decision this presentation requires?
\end{alertbox}

\subsection{Worked Example: A Board Presentation}

To make the Launchpad concrete, consider a VP of Product at a mid-size SaaS company presenting to the board. The flagship product is losing market share.

\textbf{SCQA:} S --- ``We entered 2025 with 34\% market share and 92\% renewal rates.'' C --- ``Three competitors launched AI-native alternatives; renewals dropped to 84\%; two top-20 accounts are in vendor evaluations.'' Q --- ``How do we reverse market share decline before it becomes structural?'' A --- ``We must ship an AI-native rebuild by Q3 or risk losing our SMB position permanently.''

\textbf{Big Idea:} ``Investing \$15M in an AI-native rebuild by Q3 prevents \$40M in annual churn.'' (14 words. Arguable. Specific. Clear stakes.)

\textbf{Three Pillars:} (1) The competitive threat is accelerating --- market data, competitor timelines, win/loss analysis. (2) Our current architecture cannot compete --- technical debt assessment, feature gap, customer feedback. (3) The rebuild is feasible within budget --- engineering plan, resource requirements, ROI model. MECE check: external threat, internal cause, proposed response. No overlap. No gaps.

\textbf{Temperature:} Level 4 --- Persuade. The board needs to approve significant resource reallocation. This requires overcoming objections and securing a decision.

\textbf{Audience Profile:} Know the financials and market position; do not know technical architecture details. Fear disrupting a product generating 60\% of revenue. Can decide budget reallocation up to \$15M. Value evidence-based reasoning and clear ROI timelines.

\begin{insightbox}[THE LAUNCHPAD PAYOFF]
With this Launchpad complete, the presentation practically designs itself. The slide structure follows the three pillars. The detail level matches the audience's knowledge. The emotional register matches the temperature. The evidence addresses the audience's fears. Every design decision has a strategic rationale. This is the difference between a presentation that is \emph{assembled} and one that is \emph{engineered}.
\end{insightbox}

\subsection{Why Most People Skip This}

If the Launchpad is so critical, why do most presenters skip it? Three reasons.

\textbf{It is hard.} Defining a Big Idea in under 20 words requires more intellectual effort than writing 50 slides of bullet points. As Blaise Pascal wrote: ``I would have written a shorter letter, but I did not have the time.''\endnote{Widely attributed to Pascal (``Lettres provinciales,'' 1657). The underlying insight --- that brevity requires more effort than verbosity --- is foundational to the Launchpad philosophy.}

\textbf{It is invisible.} No one sees the Launchpad. Stakeholders see slides. Managers see progress when slide counts increase. A presenter who spends three days on the Launchpad and produces zero slides appears to have accomplished nothing. The organizational incentive structure rewards visible output over invisible strategy.

\textbf{It is uncomfortable.} The Launchpad forces you to commit to a position before you have the comfort of ``exploring the data.'' You must state your Big Idea. You must choose three pillars and exclude everything else. You must set a temperature, which means deciding what you are asking the audience to do. For many presenters, this commitment feels premature and risky.\endnote{This maps to what psychologists call ``premature closure anxiety'' --- resistance to narrowing options before all information is available. See: Gary Klein, ``Sources of Power: How People Make Decisions,'' MIT Press, 1998.}

\begin{diagnosisbox}
\textbf{The Invisible Work Paradox.} The activities with the highest leverage in presentation creation --- defining the argument, profiling the audience, calibrating the temperature --- are precisely the activities that produce no visible output. Organizations that reward slide production over strategic thinking will systematically produce bad presentations. This is not a talent problem. It is a systems problem.
\end{diagnosisbox}

The consulting industry understood this decades ago. McKinsey analysts spend their first week on any engagement not gathering data but structuring the problem --- defining the key question, building the issue tree, identifying the governing thought.\endnote{Ethan Rasiel, ``The McKinsey Way,'' McGraw-Hill, 1999. Rasiel describes the ``initial hypothesis'' approach: beginning with a hypothesis about the answer \emph{before} collecting data.} Only after this intellectual architecture is complete do they begin analysis. The same principle applies to presentations. The Launchpad is the intellectual architecture. The slides are the execution.

\begin{keybox}[THE LAUNCHPAD CHECKLIST]
Before opening any presentation software, confirm:
\begin{enumerate}[leftmargin=1.5em]
  \item \textbf{SCQA} is complete and each element passes its test question
  \item \textbf{Big Idea} is a single sentence, under 20 words, with a point of view and stakes
  \item \textbf{Three Pillars} are MECE and each has at least two pieces of supporting evidence
  \item \textbf{Temperature} is explicitly chosen and matched to the audience and objective
  \item \textbf{Audience Profile} covers knowledge, fears, decision authority, and values
\end{enumerate}
If any element is missing, you are not ready to build slides. Go back.
\end{keybox}

\pullquote{The quality of a presentation is determined in the first 45 minutes of preparation --- before any slide exists. Everything after that is execution.}{Storymakers Principle}

\begin{keybox}[TRY THIS NOW]
Write the SCQA for your next presentation in four sentences --- one per element. Then test it: would someone who reads \emph{only} these four sentences understand your argument? Could they summarise what you are recommending and why? If the four sentences do not stand alone as a coherent mini-argument, the SCQA needs work. Share it with a colleague who has no context. If they can explain your recommendation back to you, the Launchpad is solid. If they cannot, revise until they can.
\end{keybox}

\clearpage
